\chapter{User Manual}
\label{ch:manual}

\section{Environment}

\begin{lstlisting}[language=bash]
python -m venv .venv
source .venv/bin/activate
pip install -r requirements.txt
\end{lstlisting}

The Phase~7 production logs report Python~3.11.7, NumPy~2.4.6,
PyTorch~2.15.0+cpu, and NGSolve~6.2.2608.
The versions in \texttt{requirements.txt} are minimum-version constraints,
not a fully pinned lockfile. For comparison with the reported cluster
results, use the versions above and record the operating system, BLAS
backend, thread counts, and NGSolve build. The Phase~1 baseline package
captures an earlier local environment; results from that environment should
not be combined with Phase~7 production metrics without qualification.

\section{Repository layout and numerical stages}

Run commands from the repository root so Python resolves the \texttt{src}
and \texttt{tests} packages. The material model
(\texttt{src/picnn.py}), its energy derivatives
(\texttt{src/constitutive.py}), and the NGSolve integration-point bridge
(\texttt{src/ngsolve\_umat.py}) are separate modules. Phase~2 separates
periodic geometry and kinematics (\texttt{src/rve.py}), mixed matrix
equations (\texttt{src/mixed\_matrix.py}), nonlinear RVE solution
(\texttt{src/rve\_solver.py} and \texttt{src/free\_particle.py}), and
homogenization. Dataset generation, integrity auditing, and training have
their own entry points, so each stage can be checked before an expensive
production run.

The workflow has two stages. Phase~1 trains on analytical neo-Hookean data
and demonstrates pICNN inference inside an NGSolve solve. Phase~2 computes
homogenized labels from periodic particle RVEs, audits and trains on those
labels, and evaluates the surrogate in a displacement FE problem. Phase~2
does not replace the small Phase~1 proof that the PyTorch--NGSolve interface
works.

\section{Phase 1 verification}

\begin{lstlisting}[language=bash]
python -m tests.verify_model_interface
python -m tests.verify_reference
python -m tests.verify_tangent
python -m tests.verify_neohooke_response
python -m tests.verify_element
python -m tests.verify_patch
python -m tests.verify_fe_comparison
\end{lstlisting}

The V4 equilibrium check uses a seeded pICNN by default and does not require a
checkpoint. To repeat it with the trained Phase~1 material, pass the checkpoint
explicitly:
\begin{lstlisting}[language=bash]
python -m tests.verify_element --ckpt checkpoints/phase1_validated.pt
\end{lstlisting}
Run the checks in the listed order because each adds one layer of
complexity: API shapes, reference normalization, constitutive accuracy,
tangent consistency, nonlinear equilibrium, patch consistency, and then a
matched analytical-versus-neural FE comparison. A passing V4 solve verifies
assembly and convergence for its loading case; V5 and V6 separately check
uniform-field reproduction and agreement with the reference material.

\section{Phase 2 verification}

\begin{lstlisting}[language=bash]
python -m tests.verify_periodic_rve
python -m tests.verify_rigid_particle
python -m tests.verify_mixed_matrix
python -m tests.verify_coupled_rve
python -m tests.verify_hill_mandel
python -m tests.verify_homogenized_tangent
python -m tests.verify_free_rotation
python -m tests.verify_translation_gauge
python -m tests.verify_free_rotation_tangent
python -m tests.verify_rve_dataset
\end{lstlisting}

\section{Dataset generation}

\begin{lstlisting}[language=bash]
python -m src.generate_rve_dataset \
  --output data/rve_phase2.npz \
  --volume-fractions 0.02 0.05 0.10 \
  --n-deformations 100 \
  --seed 1729
\end{lstlisting}

This standard configuration creates 100 deformation states and evaluates
each at three volume fractions, for 300 records. Other controls include
\texttt{--maxh}, \texttt{--max-strain},
\texttt{--determinant-floor}, \texttt{--max-condition-number},
\texttt{--bulk-modulus}, and \texttt{--microscopic-jacobian-floor}.
Compression and expansion tails are configured through
\texttt{DatasetConfig} or the parallel campaign entry point. Pass
\texttt{--work-directory} to choose a separate checkpoint directory. The
generator resumes valid completed samples by default;
\texttt{--no-resume} starts a new work set.

The adjacent JSON manifest must be retained with the NPZ. Interrupted runs
should resume from atomic per-sample checkpoints rather than overwrite
completed states.

Before training or publishing the array file, validate its manifest, hashes,
shapes, split grouping, reference states, residuals, tangent symmetry,
rotation orthogonality, and Jacobian range:
\begin{lstlisting}[language=bash]
python -m src.audit_rve_dataset --dataset data/rve_phase2.npz \
  --manifest data/rve_phase2.manifest.json
\end{lstlisting}
This audit does not rerun finite-element calculations. It checks integrity
of stored records but cannot establish physical representativeness of the
chosen RVE or mesh convergence.

\section{Surrogate and macro-FE execution}

The final fine-tuning command was:
\begin{lstlisting}[language=bash]
python -u -m src.train_rve \
  --dataset data/rve_phase6_phase7_training.npz \
  --manifest data/rve_phase6_phase7_training.manifest.json \
  --checkpoint training/partial_icnn_phase6_phase7_finetuned.pt \
  --report training/training_report.json \
  --initialize-from partial_icnn_phase6_targeted.pt \
  --epochs 3000 --batch-size 64 --learning-rate 1.0e-4 \
  --stress-weight 1.0 --tangent-weight 0.10 \
  --tangent-batch-size 8 --tangent-validation-states 12 \
  --patience 20 --validation-interval 25 \
  --seed 4103 --device cpu --torch-threads 12
\end{lstlisting}

The independent Phase~7 data were generated with
\texttt{generate\_rve\_dataset\_parallel}: 200 deformation states in total,
including 50 targeted compression states in \(J\in[0.51,0.60]\) and 50
expansion states in \(J\in[1.50,1.80]\), split fractions
\(0.05/0.05/0.90\), seed 4103, \(h=0.22\), displacement order 2,
integration order 10, and 12 workers. The archived script is
\texttt{rve\_phase7\_independent\_dataset.pbs}.

The completed macroscopic command was:
\begin{lstlisting}[language=bash]
python -u -m benchmarks.phase6_macro_fe_demo \
  --checkpoint partial_icnn_phase6_phase7_finetuned.pt \
  --output report.json --volume-fraction 0.05 --stretch 0.05 \
  --maxh 0.25 --order 1 --integration-order 4 \
  --max-iter 200 --load-steps 100 --tolerance 1.0e-8 \
  --merit-function energy --load-predictor secant
\end{lstlisting}
For the mesh study, only \texttt{--maxh} and cluster thread allocation changed
to 0.20/40 and 0.16/40. Exact PBS scripts and standard output must be retained
with each \texttt{report.json}. Failed runs must not be overwritten; their
reports document the failed load step and residual.

The training report records split counts, input normalization, epoch history,
the selected validation checkpoint, and test metrics. A fine-tuned model
keeps the parent checkpoint's feature and context normalization fixed; do
not refit that normalization on held-out states. The macro-FE command uses
the selected checkpoint, a fixed volume-fraction context, target stretch,
mesh size, element order, quadrature order, and load-step strategy. Its JSON
result is the primary numerical record and should be archived with the
matching source revision, PBS script, job output, environment versions, and
checkpoint hash.

\section{Cluster execution and output interpretation}

The supplied PBS scripts document production resource requests and the
sequence of RVE generation, model training, and macro-FE checks. Set
\texttt{NGS\_NUM\_THREADS} to the allocated NGSolve thread count and
\texttt{TORCH\_NUM\_THREADS} to the PyTorch intra-op count before launching
the job. For process-parallel dataset generation, each worker constructs its
own NGSolve mesh; mesh and solver objects are not shared across processes.
Use persistent scratch storage for the work directory, retain per-sample
checkpoints until the final NPZ and manifest pass the audit, and copy the
frozen artifacts to durable project storage.

Read macro-FE reports as solver records, not as experimental validation.
Final residuals establish equilibrium only for the saved mesh, element
order, quadrature, boundary conditions, and loading path. Compare response
quantities such as reactions or stored energy across refined meshes before
claiming convergence. Preserve failed-run reports too: failure at a
particular load increment helps distinguish load-path instability from a
broken constitutive call and should not be hidden by a later run with
different controls.

\section{Expected artifacts}

At minimum, archive the input dataset and manifest, any resumed work
directory required to reconstruct it, training report and checkpoint,
macro-FE JSON report, job script and standard output, and environment
record. Compute SHA-256 hashes for large frozen data and checkpoints. The
manifest's configuration and source hashes identify the solver inputs; a
Git commit and log should also be retained when the source repository is
available. Keep exploratory, failed, and production runs under distinct
names so each report table can be traced to the exact configuration that
produced it.

\section{Regenerating report figures}

The included plotting script recreates the verification and pilot-dataset
figures from their archived numerical sources:
\begin{lstlisting}[language=bash]
python figure_sources/generate_report_figures.py
\end{lstlisting}
The script writes vector PDF figures for LaTeX and high-resolution PNG
previews into \texttt{figures/}. Its inputs include the Phase~2 pilot NPZ and
manifest in \texttt{figure\_sources/}; the Phase~1 tangent and Newton values
are read from the archived verification logs recorded in this report. The
figures distinguish a verification diagnostic from a constitutive response
curve: in particular, the pilot plot of \(\det\Fbar\) shows loading coverage,
not a stress--strain relation, because the sampled gradients include
different distortional components.

For further result figures, retain predictions and reference targets at the
same sample IDs. A useful surrogate assessment is a parity plot for energy
and stress, supplemented by relative tangent-error distributions grouped by
volume fraction and determinant range. A component-scale mesh study should
export reaction force or stored energy at every accepted load step for each
mesh; final residuals alone cannot support a response-convergence plot. A
periodic RVE mesh rendering and a schematic of opposite-face identifications
and the rigid cavity trace would also clarify the boundary conditions. These
figures require the raw Phase~6/7 evaluation predictions and macro-FE
response histories, which are not present in this report package and must
not be reconstructed from summary metrics alone.
